\subsection{43-State Overview}

Our comprehensive analysis encompasses all 43 multi-district states, spanning minority percentages from 9.8\% (Maine) to 78.4\% (Hawaii). Table~\ref{tab:demographics_summary} summarizes the distribution across threshold categories.

\begin{table}[h]
\centering
\begin{tabular}{lccc}
\toprule
\textbf{Category} & \textbf{States (n)} & \textbf{Minority Range} & \textbf{Success Rate} \\
\midrule
High minority ($\geq$50\%)    & 6  & 52.8--78.4\% & 100.0\% \\
Above threshold (42--50\%)    & 7  & 42.4--49.9\% & 62.1\%  \\
Borderline (37--42\%)         & 5  & 36.8--41.7\% & 1.3\%   \\
Below threshold ($<$37\%)     & 25 & 9.8--36.2\%  & 0.0\%   \\
\bottomrule
\end{tabular}
\caption{Distribution of 43 states across threshold categories. Success rate = percentage of configurations achieving proportional MM target.}
\label{tab:demographics_summary}
\end{table}

Representative states from each category demonstrate the threshold pattern (see Table~\ref{tab:representative_states}).

\begin{table}[h]
\centering
\begin{tabular}{lcccc}
\toprule
\textbf{State} & \textbf{Minority \%} & \textbf{Target MM} & \textbf{Achieved MM} & \textbf{\% of Target} \\
\midrule
\multicolumn{5}{c}{\textit{Above Threshold ($\geq$42\%)}} \\
California    & 65.3\% & 34 & 44 & 129\% \\
Georgia       & 49.9\% & 7  & 8  & 114\% \\
New York      & 47.5\% & 12 & 12 & 100\% \\
Mississippi   & 44.6\% & 2  & 2  & 100\% \\
\midrule
\multicolumn{5}{c}{\textit{Borderline (37--42\%)}} \\
Illinois      & 41.7\% & 7  & 5  & 71\% \\
N. Carolina   & 39.5\% & 6  & 3  & 50\% \\
\midrule
\multicolumn{5}{c}{\textit{Below Threshold ($<$37\%)}} \\
Alabama       & 36.9\% & 3  & 2  & 67\% \\
Connecticut   & 36.8\% & 2  & 1  & 50\% \\
Pennsylvania  & 26.5\% & 5  & 0  & 0\% \\
\bottomrule
\end{tabular}
\caption{Representative states showing threshold pattern. Above 42\%: achieve 80-114\% of target. Below 42\%: achieve 0-67\% of target.}
\label{tab:representative_states}
\end{table}

\subsection{The 42\% Effectiveness Threshold}

Figure~\ref{fig:threshold} demonstrates clear threshold behavior across 43 states: states with $\geq$42\% minority achieve high effectiveness (62\% average success rate, 92\% achieve target), while states below 37\% show zero success (0\% success rate, 0\% achieve target).

\begin{figure}[h]
\centering
\includegraphics[width=0.9\textwidth]{results/figure1_50state_threshold.pdf}
\caption{Edge-weighted optimization effectiveness across 43 states. Clear demarcation at 42\% effectiveness threshold. States ordered by minority percentage (left to right). Green bars indicate states achieving proportional target; red bars indicate failure to achieve target.}
\label{fig:threshold}
\end{figure}

Figure~\ref{fig:correlation} shows the strong relationship between state minority percentage and success rate (r=0.78, p$<$1e-08), validating the threshold pattern across diverse geographic contexts.

\begin{figure}[h]
\centering
\includegraphics[width=0.9\textwidth]{results/figure2_50state_correlation.pdf}
\caption{State minority percentage vs. success rate correlation. Circles indicate states achieving proportional target (green) or failing to achieve target (red). Linear fit shows r=0.78 with p$<$0.001. Vertical line marks 42\% threshold with shaded 95\% confidence region.}
\label{fig:correlation}
\end{figure}

\subsection{Statistical Validation}

Table~\ref{tab:statistical_tests} presents comprehensive statistical validation of the 42\% threshold.

\begin{table}[h]
\centering
\begin{tabular}{lcc}
\toprule
\textbf{Test} & \textbf{Value} & \textbf{p-value} \\
\midrule
Pearson correlation (r)         & 0.7776 & $<$1e-08 \\
Spearman correlation (rho)      & 0.7848 & $<$1e-08 \\
$R^2$ (variance explained)      & 0.6047 & --- \\
\midrule
t-test (above vs below)         & t=7.39 & $<$1e-08 \\
Mann-Whitney U (non-parametric) & --- & $<$1e-08 \\
\bottomrule
\end{tabular}
\caption{Statistical validation of 42\% threshold. All tests show highly significant relationships (p$<$1e-08). Above group: minority $\geq$42\% (N=13). Below group: minority $<$37\% (N=25).}
\label{tab:statistical_tests}
\end{table}

\subsection{Practical Implications}

Based on comprehensive 43-state analysis, we establish empirically validated feasibility guidelines:

\begin{table}[h]
\centering
\begin{tabular}{lcp{7cm}}
\toprule
\textbf{Minority \%} & \textbf{Effectiveness} & \textbf{Practical Guidance} \\
\midrule
$\geq$ 50\%   & Very High (100\%)     & Proportional MM representation readily achievable; states should create proportional districts \\[0.2cm]
42--50\%      & High (62\%)           & Near-proportional representation feasible; edge-weighted optimization effective \\[0.2cm]
37--42\%      & Borderline (1\%)      & Case-specific analysis required; may achieve partial representation (50-70\% of target) \\[0.2cm]
$<$ 37\%      & Low (0\%)             & Proportional representation unlikely; geographic constraints limit MM district creation \\
\bottomrule
\end{tabular}
\caption{Feasibility guidelines for proportional MM representation based on 43-state empirical analysis. Percentages indicate average success rate (achieving proportional MM target) across all configurations within category.}
\label{tab:practical_guidelines}
\end{table}

\textbf{Legal and practical applications:}
\begin{itemize}
    \item \textbf{Courts}: Use 42\% threshold as contextual guidance for evaluating Gingles Prong 1 (geographic compactness) claims
    \item \textbf{Legislatures}: Focus MM district creation efforts on states above 42\% where proportional representation is demonstrably achievable through optimization
    \item \textbf{Plaintiffs}: Geographic feasibility arguments have strong empirical support in states $\geq$42\%; borderline states (37-42\%) require detailed case-specific analysis
    \item \textbf{Redistricting commissions}: Use edge-weighted optimization for states near threshold to maximize feasible MM district creation
\end{itemize}
